% \begin{savequote}
% Alles in Ordnung, auch wenn Sie in Panik geraten sind und
% sich beschissen fühlen.
% Ich kenne das. Es ist furchtbar, und es wird nicht das
% letzte Mal sein.
% Aber vergessen Sie nicht, was Sie gesehen haben. Regel Nr.
% 3: Bei Herzstillstand zuerst den \EE{eigenen} Puls fühlen.
% 
% % \begin{flushright}
% \qauthor{Dickie, in: Samuel Shem: \emph{House of God.} \cite{HouseOfGodDe}}
% % \end{flushright}
% 
% \includegraphics[width=4cm]{./images/wikipedia-pd/AED_Symbol}
% \end{savequote}
\def\ChapterCode{BLS}
\chapter
[\CO{[BLS]} Basic Life Support]
{Basic Life Support für professionelle Helfer}
% \AasRsN{RS VIII}
\label{chp:AED}\label{chp:BLS}
\label{topic:atemstillstand-aed}
\label{topic:kreislaufstillstand-aed}
\label{topic:atemstillstand-bls}
\label{topic:kreislaufstillstand-bls}
\label{topic:pediatric-life-support}

\ChapterIntro
{Die Kardiopulmonale Reanimation, bestehend aus
Herdruckmassage, Beatmung und Elektrotherapie mittels eines
Defibrillators, ist die Standardtherapie bei
Kreislaufstillstand. Die korrekte
und routinierte Durchführung der Reanimation hat großen
Einfluss auf das Überleben des Patienten bzw. seine
Lebensqualität nach der Entlassung aus der
Spitalsbehandlung. Die Überlebensrate bis zur
Krankenhausentlassung nach rettungsdienstlichen
Reanimationen liegt bei ungefähr 10\PC*
\cite{Erc2010Sek01De}. Die Reanimation ist die letzte
Chance für den Patienten und muss daher vom medizinischen
Fachpersonal beherrscht werden.

Dieses Kapitel behandelt die Basisreanimation mittels
eines halbautomatischen Defibrillators basierend auf den
aktuellen Leitlinien des European Resuscitation Council
(ERC). Diese ERC-Leitlinien werden ca. alle fünf Jahre den
neuesten wissenschaftlichen Erkenntnissen angepasst.}{%
\begin{description}
  \item[Maintainer:] Josef Emhofer
  \item[Autoren:] Diverse
  \item[Reviewer:] Standardreviewprozess
  \item[Version:] {\VarDocVersionTypeName} ({\VarDocVersionTypeDescription})
  \item[SHA1:] {\DocGitCommit}
\end{description}

{\TxTeilkapitelAASS}
}

%\SectionUnderConstruction{}
\clearpage
\Dictum[Dickie, in: Samuel Shem: \emph{House of God.}
\cite{HouseOfGodDe}]{Alles in Ordnung, auch wenn Sie in Panik geraten sind und
sich beschissen fühlen.
Ich kenne das. Es ist furchtbar, und es wird nicht das
letzte Mal sein.
Aber vergessen Sie nicht, was Sie gesehen haben. Regel Nr.
3: Bei Herzstillstand zuerst den \EE{eigenen} Puls fühlen.
}

\section{\HeadingKeyword{Einleitung}}

\TextSlide{\TxBeschreibung}
{%
Nach der Kontrolle der Vitalfunktionen kann sich
herausstellen, dass der Patient \EE{kein Bewusstsein} und
\EE{keine normale Atmung} hat. 
%
Daraus wird geschlossen, dass auch \EE{kein Kreislauf
vorhanden} ist. 
%
In dieser Situation liegt die Notfalldiagnose
\EEE{Atem- und Kreislaufstillstand} vor. 
Dies macht eine \EE{Herz-Lungen-Wiederbelebung} (\EE{Reanimation}) erforderlich. 
%
Dabei versucht der Helfer den Kreislauf des Patienten \enquote*{von
außen} so gut wie möglich aufrecht zu
erhalten. 
%
Dies beinhaltet einerseits eine \E{Beatmung}, 
um dem Körper Sauerstoff zuzuführen 
und andererseits die \E{Herzdruckmassage}, 
um die mechanische Herztätigkeit zu simulieren. 
%
Zusätzlich wird versucht, mittels 
\E{Elektroschocks} (Defibrillation) 
die herzeigene elektrische Herztätigkeit wiederherzustellen.
%
Beatmung, Herzdruckmassage und Defibrillation
bilden die \EE{drei Säulen der Reanimation}.
%
}{
\begin{itemize}
  \item Kein Bewusstsein
  \item Keine normale Atmung
  \item Kein Kreislauf
  \item Herz-Lungen-Wiederbelebung (Reanimation)
\end{itemize}
}{}{}{}

\TextSlide{Anmerkung: Rettungskette und Notruf}
{Entgegen der Vorgehensweise der Rettungskette wird bei
einem Atem- und Kreislaufstillstand grundsätzlich bereits
\E{nach der Kontrolle der Vitalfunktionen} der \E{Notruf
abgesetzt} und erst danach mit der \E{lebensrettenden
Sofortmaßnahme} begonnen. Die genaue Vorgehensweise und die
Techniken werden im Folgenden detailliert beschrieben.
\bigskip }{
\begin{itemize}
  \item Nach der Kontrolle der Vitalfunktionen Notruf absetzen.
  \item Anschließend lebensrettende Sofortmaßnahmen einleiten.
\end{itemize}
}{}{}{}

% \MassnahmenNG{Erste-Hilfe-Maßnahmen}{Atem- und Kreislaufstillstand}{}{
% \begin{itemize}
%   \item Herz-Lungen-Wiederbelebung (Reanimation)
%   \begin{itemize}
%     \item Herzdruckmassage
%     \item Beatmung
%   \end{itemize}
% \end{itemize}
% }

% \TextSlide{Lebensrettende Sofortmaßnahme
% }{
% Nach dem Absetzen des Notrufes muss sofort mit der
% \EE{Herz-Lungen-Wiederbelebung (Reanimation)} begonnen werden. Diese dient der
% Aufrechterhaltung einer von sich aus ausbleibenden Atmung und fehlenden Kreislauftätigkeit. 
% Zusätzlich zielt sie auf ein Wiederingangsetzen der eigenständigen
% Herz- sowie Lungentätigkeit ab. Die Reanimation besteht aus zwei
% Teilschritten, der \EE{Herzdruckmassage} zur Nachahmung einer
% Herztätigkeit und der \EE{Beatmung} zur Zuführung von Sauerstoff in die
% Lungen.
% }{
% \begin{itemize}
%     \item Herz-Lungen-Wiederbelebung (Reanimation)
%     \begin{itemize}
%         \item Herzdruckmassage
%         \item Beatmung
%     \end{itemize}
% \end{itemize} 
%  }

\section{\HeadingKeyword{Säulen} der Reanimation}



\TextSlide{3 Säulen der Reanimation}
{Die Maßnahmen der (Basis-)Reanimation setzen sich im
Wesentlichen aus drei Säulen zusammen:
\begin{enumerate}
  \item Herzdruckmassage (HDM)
  \item Beatmung
  \item Elektrotherapie (Defibrillation)
\end{enumerate} }{
\begin{itemize}
  \item Herzdruckmassage
  \item Beatmung
  \item Defibrillation
\end{itemize}
}{}{}{}

\begin{PictureSeries}{}{Die drei Säulen der Basiseanimation}
\PictureSeriesRow{3}
{./images/hirtler-aass/Reanimation-Position-2}
{Herzdruckmassage \SourcePicture{Hirtler}\label{fig:druckpunkt}}
{./images/hirtler-aass/Beutelbeatmung1}
{Beatmung \SourcePicture{Hirtler}}
{./images/hirtler-aass/Reanimation-Defi-Achtung}
{Defibrillation \SourcePicture{Hirtler}}
{}
{}
\end{PictureSeries}

\subsection{Die \HeadingKeyword{Herzdruckmassage} simuliert die mechanische Herzaktion und erzeugt einen Minimalkreislauf}
\index{Herzdruckmassage}

\Text{{\TxBeschreibung}}
{%
\DefinitionInPara{%
Bei der \T{Herzdruckmassage} (\T{HDM}) wird das Herz von außen rhythmisch durch einen Helfer komprimiert, 
jeweils gefolgt von einer Entlastungsphase. 
%
Dadurch wird die natürliche mechanische Herzaktion simuliert
und ein künstliches Fließgeschehen (Zirkulation) des Blutes erzeugt.
}%DefinitionInPara
%
Die frühzeitige Durchführung 
und das Beherrschen der richtigen Technik 
hat wesentlichen Einfluss auf das Überleben des Patienten. 
%
%Sie ist daher ausführlich beschrieben 
%und muss ausreichend und regelmäßig trainiert werden. 
}

% Die Herzdruckmassage bewirkt mittels Kompression des Herzens
% zwischen Wirbelsäule und Brustbein ein künstliches Fließgeschehen
% (Zirkulation) des Blutes.
\TextSlide{Technik}
{%
Die folgende Beschreibung orientiert sich an den Empfehlungen des \I{European Resuscitation Council} (\I{ERC}) und ist in der Bilderserie \FullPrettyRef{fig:hdm} illustriert:
%\begin{itemize}
%  \item \StepHeading{Position und Druckpunkt}: 
%  Der Patient muss auf einem
%  harten Untergrund liegen. 
%  %
%  Ggfs. ist der Patient aus einem Bett
%  o.\,ä. auf den Boden zu bringen. 
%  %
%  Der Helfer kniet dann seitlich neben dem Körper des Patienten 
%  und sucht   den richtigen \E{Druckpunkt} für die Herzdruckmassage   auf. 
%  %
%  Dieser liegt in der \E{Mitte des Brustkorbes}, 
%  d.\,h. in der unteren Hälfte des Brustbeins. 
%  %
%  Dort wird der Handballen einer Hand aufgelegt. 
%  Hebt der Helfer die Hand (unabsichtlich) vom Körper des Patienten ab, 
%  so muss der Druckpunkt neu aufgesucht werden, 
%  bevor die nächste Herzmassage anschließt.
%  %
%  %
%  \item \StepHeading{Arm- und Körperhaltung}: 
%  Die richtige Arm- und Körperhaltung ist wichtig, 
%  um mit wenig Kraftaufwand die bestmögliche Leistung zu erzielen.
%  Nachdem der \E{Handballen} einer Hand \E{auf den Druckpunkt} gelegt wurde, 
%  wird der Handballen der zweiten Hand direkt darüber gelegt. 
%  %
%  Die \E{Finger werden ineinander verschränkt}. 
%  %
%  Nun werden die \E{Arme durchgestreckt} 
%  und der Oberkörper solange nach vorne gebeugt bis sich die
%  \E{Schultern genau über dem Druckpunkt} befinden. 
%  %
%  Dadurch wird der Druck direkt auf das darunter liegende Brustbein 
%  und nicht seitlich auf die Rippen ausgeübt. 
%  %
%  Für die \E{Kompression} des Herzens drückt man mit Hilfe des eigenen Gewichts direkt nach unten.
%  %
%  Das Herz wird so zwischen Brustbein und Wirbelsäule zusammen  gedrückt 
%  und Blut wird in die Aterien ausgeworfen.
%  %
%  %
%  \item \StepHeading{Druckerzeugung und Druckentlastung}: 
%  Der \E{Druck} wird
%  -- wie bereits erwähnt -- 
%  \E{direkt nach unten} ausgeübt.
%  %
%  Der Druck hat \E{kräftig, aber nicht ruckartig} zu erfolgen.
%  %
%  Die \E{Eindrücktiefe} beträgt \E{\VonBis{5}{6}\CM*}, 
%  beim durchschnittlichen Erwachsenen entspricht das ungefähr   \nicefrac{1}{3} des Brustkorbes. 
%  %
%  Der \E{Brustkorb} muss sich nach der Kompression wieder \E{vollständig entlasten} können
%  damit das Herz wieder mit Blut befüllt wird. 
%  %
%  Während der Entspannungsphase darf der Helfer daher nicht auf dem   Brustkorb des Patienten lehnen. 
%  %
%  Die \E{Druck- und Entlastungsphasen} dauern beide \E{gleich lange}
%  \cite{Erc2010Sec2}.
%  
%  % \EE{Druckentlastung} Der Brustkorb muss sich wieder ganz
%  % ausdehnen können. Es darf daher \E{kein Restdruck} bestehen. Die Druck-
%  % und Entlastungsphasen dauern beide gleich lange.
%  
%  \item \StepHeading{Druckfrequenz}: 
%  Die Druckfrequenz
%  (\enquote{Arbeitsgeschwindigkeit}) der Herzdruckmassage
%  beträgt \E{mindestens 100 Massagen pro Minute} (aber nicht mehr als
%  120 Massagen pro Minute).
%  
%  \item \StepHeading{Verhältnis}: 
%  Die Herzdruckmassage wechselt sich mit der Beatmung des Patienten ab.
%  % 
%  Der seitlich am Patient knieende Helfer beginnt die Wiederbelebung mit \E{30 Herzdruckmassagen}, 
%  darauf folgen \E{zwei Beatmungen} durch den am Kopf befindlichen Helfer,
%  danach wird sofort mit 30 HDM fortgefahren usw. 
%  %
%  \E{Erfolglose Beatmungsversuche} werden nicht wiederholt 
%  um die HDM nicht zu verzögern. 
%  %
%  Das Wechselspiel aus 30 HDM und zwei Beatmungen wird fortgeführt
%  bis der AED einsatzbereit ist. 
%  %
%  Dies soll so rasch wie möglich geschehen. 
%\end{itemize}
%
Nachdem der Patient auf eine harte Unterlage (ggf. auf den Boden) gebracht wurde, muss der Brustkorb frei gemacht und der \EE{Druckpunkt} für die Herzdruckmassage aufgesucht werden. Dieser liegt in der Mitte des Brustkorbs, d.\,h. in der unteres Hälfte des Brustbeins. Auf diesem Punkt wird ein Handballen des Helfers gelegt. Der Handballen der anderen Hand kommt auf die erste Hand, direkt über dem Druckpunkt. Die Finger werden ineinander verschränkt. Der für die Kompression notwendige Druck wird erzeugt, indem sich der Helfer über den Patient beugt, die Arme durchstreckt und mit Hilfe seines eigenen Gewichts gerade nach unten drückt. Dabei soll eine \EE{Eindrücktiefe} von  \E{\VonBis{5}{6}\CM*} erreicht werden.\footnote{Beim durchschnittlichen Erwachsenen entspricht das ungefähr \nicefrac{1}{3} des Brustkorbes.} Nach der Kompression folgt die vollständige Entlastung des Brustkorbs ehe eine neuerliche Kompression beginnt.\footnote{Vollständige Entlastung bedeutet, dass der Helfer keinerlei Druck mehr auf den Brustkorb des Patienten ausübt. Die Handballen werden also in ihre Ausgangslage zurückgeführt, wobei sie gerade nicht vom Brustkorb abheben. Sollte es versehentlich passieren, dass der Helfer den Kontakt zum Patient verliert, muss der Druckpunkt neu aufgesucht werden.} Die Abfolge von Druck- und Entlastungsphasen soll \E{harmonisch} erfolgen, d.\,h. dass beide Phasen gleich lange dauern und nicht ruckartig ablaufen \cite{Erc2010Sec2}. Die \EE{Druckfrequenz} (\enquote{Arbeitsgeschwindigkeit}) der Herzdruckmassage beträgt \E{mindestens 100 Massagen pro Minute}.\footnote{Die Herzdruckmassage soll nicht schneller als mit 120 Massagen pro Minute durchgeführt werden.}

}{
\begin{itemize}
  \item Druckpunkt: \E{Mitte} des Brustkorbes
  \item Arm- und Körperhaltung
  \begin{itemize}
    \item Handballen am Brustbein,
    \item Finger verschränkt,
    \item Arme durchgestreckt,
    \item Schultern genau über dem Druckpunkt.
  \end{itemize}
  \item Druckerzeugung und Druckentlastung
  \begin{itemize}
    \item kräftig aber nicht ruckartig
    \item Eindrücktiefe: 5-6\CM*
    \item vollständige Entlastung
  \end{itemize}
  \item Druckfrequenz: mind. 100\PerMin*
\end{itemize}
}{}{}{}

\begin{PictureSeries}{}{Bilderserie: \EE{Herzdruckmassage} \SourcePicture{Hirtler}\label{fig:hdm}}
\PictureSeriesRow{3}
{./images/hirtler-aass/Reanimation-Druckpunkt}
{Druckpunkt: Mitte des Brustkorbs. Die Finger werden
ineinander verschränkt, \ldots}
{./images/hirtler-aass/Reanimation-Position-1}
{{\ldots} die Arme durchgestreckt \ldots}
{./images/hirtler-aass/Reanimation-Position-2}
{{\ldots} und Oberkörper nach vorne gebeugt:
$\rightarrow$~Druck direkt auf Brustbein}
{}
{}
\end{PictureSeries}

\Text{HDM: sooft wie möglich!}{%
Die Herzdruckmassage muss so rasch wie möglich beginnen und
so unterbrechungsfrei wie möglich durchgeführt werden.
Nach Abgabe des Schocks muss die Herzdruckmassage
sofort weitergeführt werden! Auch wenn die Defibrillation
erfolgreich war, d.\,h. ein ordnungsgemäßer Herzrhythmus
wiederhergestellt wurde, kann es einige Zeit dauern, bis
wieder ein ausreichender Kreislauf einsetzt
\cite{Sandroni2008,Skhirtladze2010}.
}{}{}{}{}

\begin{Synopsis}
  \SynopsisItem{Die Unterbrechung der HDM durch den AED/Defibrillator
  soll insgesamt nicht mehr als \E{5\Sek*} betragen!}
\end{Synopsis}


\TextSlide{Mögliche Fehlerquellen}
{
%Der größte Fehler, den ein potenzieller Helfer machen
%kann, ist, \EE{nichts zu tun}. Die professionelle
%Reanimation ist die letzte Chance des Patienten. 
%
%Oft wird vergessen, dass der Patient bei der
%Herzdruckmassage auf einer \EE{festen Unterlage} liegen
%muss. Dies kann besonders dann passieren, wenn der Patient
%im Bett vorgefunden wird, ein \E{Bett ist jedoch keine feste
%Unterlage}.
%
%Während der Herzdruckmassage kann es zu \EE{Verletzungen
%durch den Helfer} kommen.
%Dazu zählt der Bruch des Brustbeines durch zu festen Druck,
%Lungen- oder Herzquetschungen und Rippenbrüche durch
%Abweichung des Druckpunktes, seitlich ausgeübten Druck oder
%durch einen zu tiefen Druckpunkt. Es gilt jedoch: Gebrochene
%Rippen und gequetschte Organe können wieder heilen. Das
%Herz, das nicht mehr schlägt, geht vor!
%
%Ein zu ruckartiger Druck führt zu einer zu kurzen
%Kompression des Herzens, was wiederum eine verkürzte
%Auswurfphase bedeutet. Dies führt zu einer \EE{ungenügenden
%Zirkulation} des Blutkreislaufs.
Eine effiziente Herzdruckmassage erfordert einen flüssigen Ablauf 
sowie die Vermeidung von Fehlern. 
Aufgrund der Wichtigkeit der Wiederbelebung ist an dieser Stelle ausdrücklich auf mögliche Fehlerquellen hingewiesen:
\begin{itemize}
	\item \EE{Leerlauf}: Der größte Fehler ist, nichts zu tun. 
	Das kann insbesondere nach der Defibrillation oder nach den Beatmungsversuchen passieren.
	\item Ein \E{Bett ist keine feste Unterlage}. 
	Es darf nicht vergessen werden, den Patient auf eine \EE{harte Unterlage} zu legen.
	\item Wird der Druck nicht direkt von oben und/oder zu kräftig ausgeübt, 
	kann es zu \EE{Verletzungen durch den Helfer} kommen. 
	Es ist daher besonders darauf zu achten, 
	dass sich die Schultern des Helfers direkt über dem Brustbein des Patienten befinden. 
	Besteht der Verdacht den Patient verletzt zu haben, 
	muss die Herzdruckmassage trotzdem fortgesetzt werden.
	\item Ein zu ruckartiger Druck führt zu einer zu kurzen Auswurfphase des Herzens 
	und somit zu einer \EE{ungenügenden Zirkulation}.
\end{itemize}
}{
\begin{itemize}
  \item Nichts tun
  \item Keine feste Unterlage
  \item Verletzungen durch falschen Druckpunkt oder falsche
  Druckrichtung
  \item Ungenügende Zirkulation durch zu ruckartigen Druck
\end{itemize}
}{}{}{}

\begin{Synopsis}
  \SynopsisItem{Wird bei der Herzdruckmassage ein falscher
  Druck ausgeübt oder eine Rippe gebrochen: Nicht aufhören,
  sondern den richtigen Druckpunkt erneut aufsuchen und die
  Herzdruckmassage fortsetzen!}
\end{Synopsis}


\subsection{\HeadingKeyword{Beatmung}}
\index{Beatmung!BLS}

Durch die Beatmung wird dem Körper Sauerstoff zugeführt und
Kohlendioxid abgeführt. Dadurch kann ein minimaler
Sauerstoffaustausch künstlich aufrecht erhalten werden.

\TextSlide{Arten und Hilfsmittel}{
Während in der ersten Hilfe die Beatmung mittels
Mund-zu-Mund-, Mund-zu-Nasen- und
Mund-zu-Mund/Nasen-Beatmung durchgeführt wird, beatmet
medizinisches Fachpersonal mit Beatmungsbeutel,
Sauerstoffreservoir und Sauerstoffzufuhr.
\VARIANT{VAR-STANDARD}{
Diese
Beatmungshilfen sind in Kapitel \FullPrettyRef{chp:ASB}, beschrieben.
Bei spezieller Ausbildung/Freigabe kann der professionelle
Helfer zum Freihalten der Atemwege einen Guedel-Tubus
verwenden.
}
}{
  \begin{itemize}
    \item Beatmungsbeutel
    \item Sauerstoffzufuhr, Reservoir
  \end{itemize}
}{}{}{}







\TextSlide{Technik und Parameter der Atmung}
{
Das Beherrschen der richtigen Technik ist gerade bei der
Beatmung sehr schwierig. Regelmäßiges Training ist daher
unbedingt notwendig. Die folgende Beschreibung orientiert
sich an den Empfehlungen des European Resuscitation Council
(ERC).
\VARIANT{VAR-STANDARD}{ Der \EE{Umgang mit Sauerstoff und
Beatmungshilfen} (Sicherheitshinweise, C-Griff,
Technik,\ldots) ist unter \prettyref{topic:sauerstoff} und
\FullPrettyRef{topic:beatmung-hilfsmittel}, beschrieben.}

Die \EE{Atemfrequenz} des Erwachsenen beträgt
\VonBis{12}{16}\PerMin*, d.\,h. wir atmen
\VonBis{12}{16}\PerMin* ein und aus. Das
\EE{\underline{maximale} (!) Atemzugsvolumen} in {\Ml}
entspricht dem 10-fachen des Körpergewichtes, d.\,h. ein
90\Kg* schwerer Mann kann 900\Ml* Luft ein- und wieder
ausatmen.
Ein normaler Atemzug beträgt jedoch nur etwa
500\Ml*\VARIANT{VAR-STANDARD}{ (vgl.
\FullPrettyRef{topic:atemvolumina})}. Dieser Wert soll auch
bei der Beatmung angestrebt werden. Dazu wird der
Beatmungsbeutel gleichmäßig zusammen gedrückt bis eine
sichtbare Atembewegung vorhanden ist. Der
\EE{Beatmungsvorgang} soll in etwa \EE{1\Sec*} dauern
\cite{Erc2010Sek02De}. Ein Gefühl für die richtige Dosierung
kann durch ausreichend Training in der
Beutel-Masken-Beatmung erreicht werden.

Die ausreichende \ce{O2}-Konzentration in der
Umgebungs-/Einatmungsluft ist die Voraussetzung für eine
funktionierende Atmung. Ist sie nicht gegeben, ist auch die
beste Atemarbeit nutzlos. Der normale \ce{O2}-Gehalt der
Umgebungsluft beträgt 21\PC*. Der \ce{O2}-Gehalt in der
\EE{Ausatemluft} beträgt noch 16\PC*, d.\,h. der Mensch
verbraucht  nur 5\PC* des eingeatmeten Sauerstoffs. Bei
Beatmung mit Sauerstoffreservoir und Sauerstoffzufuhr von
15\LPerMin* kann die Sauerstoffkonzentration auf fast
100\PC* erhöht werden.


Bei entsprechender Ausbildung und Freigabe kann zum
Freihalten der Atemwege zusätzlich zum
Überstrecken des Kopfes ein
\II{Guedel-Tubus} verwendet werden. 
}{
\begin{itemize}
  \item Umgang mit Sauerstoff und Beatmungshilfsmittel
  \item Atemfrequenz \VonBis{12}{16}\PerMin*
  \item Atemzugvolumen 500\Ml*
  \item Dauer Beatmungsvorgang 1\Sec*
  \item Erfolglosen Beatmungsversuch nicht wiederholen
  \item Bei entsprechender Ausbildung/Freigabe
  Guedel-Tubus
\end{itemize}
}{}{}{}

\TextSlide{Mögliche Fehlerquellen}
{
Ein \EE{zu hohes Atemzugsvolumen} führt zu einem höheren
Druck im Nasen-Rachenraum, sobald die Lunge vollständig mit
Luft befüllt ist. Die Luft findet nur den Ausweg über die
Speiseröhre in den Magen. Dies kann während einer
Herz-Lungen-Wiederbelebung meist unbemerkt geschehen. Bei
einem \EE{zu niedrigem Atemzugsvolumen} ist die
\ce{O2}-Versorgung nicht sichergestellt, da die Luft nicht
bis in die Lunge kommt. Bei Bartträgern kann es zu einer
\EE{ungenügenden Abdichtung} kommen. Eine \EE{zu hohe
Atemfrequenz} hindert den Patienten am Ausatmen. Wenn es
beim Freimachen der Atemwege zu einer \EE{mangelnden
Überstreckung} des Kopfes kommt, kann es zu einer
unzureichenden Belüftung der Lunge kommen. Die Qualität der
Beatmung kann durch \EE{Beobachtung des Brustkorbs}
überprüft werden. Dazu wird nach jeder Atemspende der
Brustkorb des Patienten beobachtet. Dieser muss sich
deutlich heben und senken.
}{
\begin{itemize}
  \item zu hohes Atemzugsvolumen ${\rightarrow}$ Luft in den Magen
  \item zu niedriges Atemzugsvolumen ${\rightarrow}$ \ce{O2}-Versorgung nicht
  sichergestellt
  \item ungenügende Abdichtung bei Bartträgern
  \item mangelnde Überstreckung ${\rightarrow}$ unzureichenden Belüftung der
  Lunge
  \item Brustkorb beobachten ${\rightarrow}$ deutlich heben und senken
\end{itemize}
}{}{}{}

\subsection{Die \HeadingKeyword{Defibrillation} soll die normale elektrische Herztätigkeit wiederherstellen}

Das Herz ist die Pumpe des Blutkreislaufes und leistet eine
beachtliche mechanische Arbeit. Damit sie funktioniert wird
die Muskelarbeit mittels \EE{elektrischer Reize} gesteuert
(Reizleitungssystem,
\FullPrettyRef{topic:reizleitungssystem}). Ohne diese
elektrischen Ströme kommt es zu keiner Muskelarbeit des
Herzens und somit zum Kreislaufstillstand.
Oft ist die Störung des elektrischen Herzrhythmus die
Ursache des Kreislaufstillstands. Liegt eine solche Störung
vor, so kann man mittels \EE{Defibrillation}
(Elektrotherapie mittels Elektroschock) versuchen, einen normalen Rhythmus
wiederherzustellen.

\TextSlide{Defibrillatoren}{
Geräte für die Defibrillation werden \II{Defibrillator}
genannt. Es gibt folgende Gerätekategorien:
\begin{itemize}
  \item Halbautomatische Defibrillatoren
  (\II{AED}\footnote{\Lang{Abkz.} AED: Automatisierter
  externer Defibrillator }):
  Das Gerät kontrolliert den elektrischen Herzrhythmus und
  beurteilt, ob ein Schock abgegeben werden soll. Die
  Schockauslösung geschieht durch den Anwender. Diese
  Gerätekategorie kann von nicht-ärztlichem Personal und
  sogar auch von nicht ausgebildeten Laienhelfern verwendet
  werden.
  \item Manuelle Defibrillatoren: Das Gerät zeigt den
  elektrischen Herzrhythmus an, der Anwender muss
  entscheiden, ob ein Schock abgegeben werden soll. Manuelle
  Defibrillatoren sollen nur durch ärztliches Personal
  angewandt werden. Viele neuere Geräte verfügen über einen
  \enquote{Halbautomaten-Modus}, in welchem das Gerät als
  AED verwendet werden kann.
  
  \VARIANT{VAR-STANDARD}{Die manuelle Defibrillation wird unter
  \FullPrettyRef{chp:ALS}, besprochen.}
  \item \Z{Automatische implantierbare Defibrillatoren
  können Patienten -- ähnlich einem Herzschrittmacher --
  implantiert werden.}
\end{itemize}
}{
\begin{itemize}
  \item halbautomatische Defibrillatoren (AED)
  \item manuelle Defibrillatoren
  \item \Z{automatische implantierbare Defibrillatoren}
\end{itemize}
}{}{}{}

% Der AED wird mittels zweier Klebeelektroden mit dem
% Patienten verbunden. Mittels dieser Elektroden wird
% einerseits der vorhandene elektrische Herzrhythmus
% abgeleitet und vom Gerät beurteilt, und andererseits bei
% Bedarf der Schock abgegeben. 

\begin{PictureSeriesWide}{}{Bilderserie:
\EE{Defibrillation}\label{fig:defibrillation}}
\PictureSeriesRowWide{3}
{./images/wikipedia-cc-by-sa/Defib_electrode_placement}
{Anbringen der Klebeelektroden \SourcePicture{WMC:Andersat,
CC-BY-SA-3.0}}
{./images/hirtler-aass/Reanimation-Defi-Achtung}
{Bei der Analyse und Schockabgabe darf niemand den
Patienten berühren!}
{./images/wikipedia-cc-by-sa/Semi_automatic_defi_with_electrodes}
{Automatisierter Externer Defibrillator (AED) mit
angeschlossenen Klebeelektroden für den Einmalgebrauch
\SourcePicture{WMC:Hborkyb, CC-BY-SA-2.5-US}}
{./images/wikipedia-pd/AED_Symbol}
{Hinweisschild \enquote{Defibrillator}}
\end{PictureSeriesWide}

\TextSlide{Elektrische Herzrhythmen}{
Die Defibrillation ist nicht bei jeder Störung des
Herzrhythmus sinnvoll. Man unterscheidet daher schockbare
und nicht-schockbare Rhythmen:
\begin{itemize}
  \item \EE{Schockbare Rhythmen}
\begin{itemize}
    \item \II{Kammerflimmern}: Unkoordnierte, \enquote{wirre}
    elektrische Aktivität
    \item \II{Pulslose Ventrikuläre Tachykardie} (\I{PVT}):
    Zu schneller Herzrhythmus, das Herz kann sich nicht
    ausreichend füllen, daher fehlt die Pumpleistung
  \end{itemize}
  \item \EE{Nicht-schockbare Rhythmen}
  \begin{itemize}
    \item \II{Asystolie}: Keine Elektrische Aktivität
    \item \II{Pulslose elektrische Aktivität} (\I{PEA}): Das
    Herz reagiert nicht auf elektrische Reize
  \end{itemize}
\end{itemize}
Beim AED übernimmt das Gerät die Beurteilung ob ein
elektrischer Herzrhythmus schockbar oder nicht schockbar
ist. Die Herzrhythmen werden im Kapitel
\FullPrettyRef{topic:herzrhythmusstoerungen}, genauer
besprochen.
}{
\begin{itemize}
  \item schockbare Rhythmen
  \begin{itemize}
    \item Kammerflimmern
    \item pulslose ventrikuläre Tachykardie
    \end{itemize}
  \item nicht-schockbare Rhythmen
  \begin{itemize}
    \item Asystolie
    \item pulslose elektrische Aktivität
    \end{itemize}
\end{itemize}
}{}{}{}



\begin{figure}[h]
\caption{Sinusrhythmus und reanimationspflichtige Rhythmen}

\subfloat[\Z{(Langsamer) Sinusrhythmus}
\SourcePicture{WM/PD}] {\includegraphics[width=\linewidth]{./images/wikipedia-pd/Lead_II_rhythm_generated_sinus_bradycardia}}

\subfloat[\Z{Reanimationspflichtig, schockbar: Ventrikuläre Tachykardie} \SourcePicture{WM/PD}]
{\includegraphics[width=\linewidth]{./images/wikipedia-pd/Lead_II_rhythm_ventricular_tachycardia_Vtach_VT}}

\subfloat[\Z{Reanimationspflichtig, schockbar: Kammerflimmern}
\SourcePicture{WM/PD}]
{\includegraphics[width=\linewidth]{./images/wikipedia-pd/Lead_II_rhythm_generated_ventricular_fibrilation_VF}}

\subfloat[\Z{Reanimationspflichtig, nicht schockbar: Asystolie}
\SourcePicture{WM/PD}]
{\includegraphics[width=\linewidth]{./images/wikipedia-pd/Lead_II_rhythm_generated_asystole}}
\end{figure}



\TextSlide{Verwenden eines AED}{
Grundsätzlich ist jedes Gerät so zu verwenden wie es der
Hersteller in der \EE{Bedienungsanleitung} vorschreibt. Bei
der Verwendung von AEDs kommt es meist zu folgendem
allgemeinen Vorgehen:

Das Gerät wird in unmittelbarer Reichweite zu Patient und
Helfer positioniert, die Gerätetasche wird geöffnet. Die
Klebeelektroden sind Einmalprodukte und daher zusätzlich
verpackt. Nach Öffnen der Schutzhülle werden die Kabel der
Elektroden an das Gerät angeschlossen. Danach wird einzeln
von jeder der beiden Elektroden die Schutzfolie abgezogen
und die Elektrode \E{faltenfrei} auf den Patienten geklebt.
Ist die Schutzfolie einmal abgezogen darf die Elektrode
nicht mehr abgelegt werden, da die Klebeschicht sonst
verschmutzt würde und wirkungslos wird. Ggf. müssen die
Brusthaare des Patienten vor dem Aufkleben abrasiert werden.
Ein Einmalrasierer ist i.\,d.\,R. in der Gerätetasche
vorrätig.
Schlecht geklebte Elektroden (Falten) oder nicht
abrasiertes, dichtes Brusthaar führen zu einem unnötig hohen
elektrischen Widerstand im Stromfluss und machen die
Elektroschocks daher ineffizient. Die \EE{Positionierung der
Elektroden} ist in Abbildung \ref{fig:defibrillation}
ersichtlich. Durch diese Positionierung ist sicher gestellt,
dass ein Großteil des elektrischen Feldes zwischen den
Elektroden durch das Herz läuft.

Sind die Elektroden angeklebt, kann mit der Beurteilung des
elektrischen Herzrhythmus begonnen werden. Die
\EE{Rhythmusanalyse} (kurz Analyse) führt das Gerät auf
Knopfdruck selbständig durch. Es ertönt eine Warnung, z.\,B.
\Speech{\enquote{Analyse läuft! Patienten nicht berühren!}}
Der Patient (oder elektrisch leitende Gegenstände (z.\,B.
Metalle, Wasser) die mit dem Patient in direktem Kontakt
stehen) dürfen vom Helfer nicht berührt werden. Nach der
Analyse wird den Helfern per Sprachausgabe mitgeteilt, ob
ein schockbarer oder ein nicht-schockbarer Rhythmus
vorliegt. Wird ein Schock empfohlen, kann mittels
Tastendruck der Elektroschock ausgelöst werden.
Selbstverständlich darf auch während der \EE{Schockabgabe}
kein Kontakt zum Patienten besteht.
\EE{Helfer müssen hier sorgfältig auf Sicherheit achten
(warnen $\rightarrow$ zurücktreten $\rightarrow$
kontrollieren)!}

\subparagraph{Sicherheits- und weitere Bedienungshinweise}:
\begin{itemize}
  \item Der Defibrillator darf nicht in einer \EE{nassen
  Umgebung} oder auf \EE{Metallunterlagen} eingesetzt werden,
  da Wasser und die meisten Metalle den elektrischen Strom
  leiten.
  Dadurch sind Helfer bei der Schockabgabe gefährdet. 
  Patienten, die aus dem Wasser gerettet werden,
  müssen daher vor der Schockabgabe abgetrocknet und auf
  trockenen Untergrund gebracht werden.
  \item Viele Geräte dürfen bei \EE{Kindern} erst ab einem 
  bestimmten Körpergewicht oder Alter (je nach Herstellerangabe)
  eingesetzt werden.
  \item Defibrillatoren dürfen nicht in Umgebungen
  eingesetzt werden, in welchen \EE{Explosionsgefahr}
  herrscht.
  \item Während der Schockabgabe sind etwaige
  \E{Sauerstoffquellen} so weit wie möglich von den
  Elektroden weg zu halten (Richtwert ca. 1\Meter*).
%   \item \E{Unterkühlte Patienten} mit einer
%   Körperkerntemperatur unter 30\DegC{} dürfen ebenfalls
%   nicht defibrilliert werden. % ex mit ERC 2010
  \item Folgende Umstände können zu einem \E{Abbruch der
  Analyse} führen: Eigenbewegung des Patienten,
  Manipulationen am Patienten (durch unachtsame Helfer),
  Manipulationen am Defibrillator oder Kabel, maschinelle
  Beatmung, Transport, sonstige Erschütterungen. In diesem
  Fall kann kein Schock abgegeben werden!
  \item \I{Medikamentenpflaster} müssen vor der
  Defibrillation entfernt werden.
  \item Die Elektroden müssen \E{ohne Lufteinschlüsse}
  aufgeklebt werden.
  \item Bei Starkstromleitungen muss ein Abstand von
  mindestens 3\Meter* gehalten werden, da sonst die Analyse
  beeinflusst werden kann.
\end{itemize}
Zusätzliche Sicherheits- und Bedienhinweise sind in der
Bedienungsanleitung des jeweiligen Geräts nachzulesen! }{
\begin{itemize}
  \item Verwenden gem. Bedienungsanleitung!
  \item Positionierung
  \item Aufkleben der Elektroden
  \item Analyse
  \item Schockabgabe
  \item Sicherheitshinweise
\end{itemize}
}{}{}{}

\Text{{\TxQuerverweise}}{%
\begin{itemize}
  \item Herz: \dotfill \FullPrettyRef{topic:herz};
  \item Reizleitungssystem:
  \dotfill \FullPrettyRef{topic:reizleitungssystem}
  \item EKG: \dotfill \FullPrettyRef{topic:ekg}
  \item Herzrhythmen:
  \dotfill \FullPrettyRef{topic:herzrhythmusstoerungen}
\end{itemize}
}{}{}{}{}

\clearpage
\section{\HeadingKeyword{Algorithmus} Herz-Lungen-Wiederbelebung}

\TextSlide{Einleitung}{Der Algorithmus der Reanimation
orientiert sich an den Empfehlungen der neuesten
ERC-Leitlinien des so genannten \EE{Basic Life Support}
(\EE{BLS}). Ziel der ERC-Leitlinien ist eine einheitliche,
bestmögliche Vorgehensweise bei der Reanimation. Diese hier
beschriebenen Basismaßnahmen der Herz-Lugen-Wiederbelebung
sind durch zahlreiche Studien wissenschaftlich abgesichert
und werden  alle fünf Jahre aufgrund neuer Erkenntnisse
aktualisiert.\footnote{\Z{In Österreich gibt es den
Österreichischen Rat für Wiederbelebung (Austrian
Resucitation Council, ARC), der unter Wahrung der
wissenschaftlichen Herangehensweise und in Zusammenarbeit
von medizinischen Fachgesellschaften und
Rettungsorganisationen (ASBÖ, Johanniter, MA70, Malteser und
Rotes Kreuz) die Bereitschaft zur Reanimation und die nötige
Ausbildung fördert.}} Neben dem Basic Life Support, welcher
hauptsächlich für Laien und nicht-ärztliche Helfer ausgelegt
ist, ist der \EE{Advanced Life Support} (Erweiterte
Maßnahmen, \EE{ALS}), welcher für speziell ausgebildete
Notfallteams konzipiert ist, weit verbreitet. NEF- und
NAW-Teams arbeiten i.\,d.\,R. nach diesem Algorithmus. Für
Kinder ist der \EE{Pediatric Life Support} (\EE{PLS})
anzuwenden.
% 
% \Z{
% Es ist daher in der Ausbildung notwendig, diese Empfehlungen
% so originalgetreu wie möglich widerzugeben, weshalb sich die
% Formulierungen der folgenden Vorgehensweisen nur geringfügig
% von den Texten aus \cite{Erc2010Sek02De,Erc2010Sek03De,Erc2010Sek04De,Erc2010Sek06De} unterscheiden.}
}{
\begin{itemize}
  \item orientiert sich an den ERC-Leitlinien
  \item verschiedene Formen:
  \begin{itemize}
    \item Basic Life Support (BLS)
    \item Advanced Life Support (ALS)
    \item Pediatric Life Support (PLS)
    \end{itemize}
\end{itemize}
}{}{}{}

\subsection{Basic Life Support (BLS)}
\TextSlide{Vorgehensweise BLS}{
Die ERC betont die Wichtigkeit frühzeitigen Handelns.
Durch engmaschiges Monitoring von Notfallpatienten soll ein
drohender Kreislaufstillstand frühzeitig erkannt werden,
sodass rechtzeitig Reanimationsbereitschaft hergestellt
wird. Nur so kann im Falle eines tatsächlichen
Kreislaufstillstandes der Algorithmus sofort beginnen und
unterbrechungsfrei abgearbeitet werden. Mit jeder Minute
Verzögerung sinkt die Überlebenswahrscheinlichkeit des
Patienten um ca. 10\PC*.

Ist der Patient (bereits beim Antreffen) reglos, so geht man
wie folgt vor:
\begin{itemize}
  \item \EE{Szeneüberblick mit (Selbst-)Schutz}: Es ist auf
  Sicherheit zu achten. Helfer und Patient dürfen nicht
  gefährdet sein. \Z{Dieser Punkt ist ident mit dem
  allgemeinen Vorgehen des Patientenmanagements.}
  \item \EE{Eindruck/Hauptbeschwerde}: Der Patient ist
  reglos. Bei Atemnot/Atemstill\-stand kann Blässe und/oder
  Zyanose auffallen. Eventuell ist ein Unfallhergang
  ersichtlich.
  \Z{Dieser Punkt ist ident mit dem allgemeinen Vorgehen 
  des Patientenmanagements.}
  \item \EE{Ansprechbarkeit prüfen, Bewusstsein beurteilen}:
  Der Patient ist laut und deutlich anzusprechen. Dabei kann
  man gleichzeitig an den Schultern des Patienten rütteln
  (Ausnahme: Traumageschehen!). Ein Schmerzreiz wird im
  Bereich des Kopfes (Wange, Nasenscheidewand, Kieferwinkel)
  gesetzt.
  \item Ist der \EE{Patient ansprechbar} folgt das Vorgehen
  den allgemeinen Regeln des Patientenmanagements.
  \item Ist der \EE{Patient nicht ansprechbar} wird ein
  Notarzt angefordert. Idealerweise deligiert man diese
  Aufgabe an einen zusätzlichen Helfer um das weitere
  Vorgehen nicht zu verzögern.
  \item \EE{Kopf überstrecken}: Um die Atemwege zu öffnen
  und frei zu halten wird der Kopf überstreckt. 
%   Dabei legt
%   man eine Hand auf die Stirn und zieht den Kopf leicht 
%   nach hinten. Die
%   Fingerspitzen der anderen Hand heben dabei gleichzeitig
%   das Kinn an. Ggf. muss der Patient vorher auf den Rücken
%   gedreht werden.
  \item \EE{Atmung prüfen}: Das Ohr des Helfers wird
  knapp über den Mund des Patienten gehalten. Die
  Blickrichtung des Helfers fällt auf den Brustkorb des
  Patienten. Jetzt wird für die Dauer von 10\Sec* die
  Atmung des Patienten beurteilt:
  \begin{itemize}
    \item \E{Sehen}: Bewegungen des Brustkorbs
    \item \E{Hören}: Atemgeräusche
    \item \E{Fühlen}: Luftstrom der Ausatemluft
    \end{itemize}
  \item \EE{Beurteilung} ob die \EE{Atmung normal} ist: Das
  ERC versteht unter \enquote{normaler} Atmung jede
  Atemtätigkeit die mit selbständiger Aufrechterhaltung der
  Lebensfunktionen vereinbar ist. \E{Schnappatmung}
  (vereinzelte, langsame oder geräuschvolle Atemzüge mit
  welchen der Patient nach Luft schnappt)
  \E{ist keine normale Atmung!} Der Schnappatmung folgt oft
  Atemstillstand.
  \item Liegt eine \EE{normale Atmung} vor, wird der Patient
  in die stabile Seitenlage gedreht. Danach folgen alle
  weiteren {\TxMassVit}.
  \item Liegt \EE{keine normale Atmung} vor, muss sofort mit
  der Herz-Lung\-en-Wied\-er\-be\-leb\-ung begonnen werden. Die
  notwendigen Hilfsmittel sollen zu diesem Zeitpunkt bereits
  am Einsatzort sein. Der seitlich am Patient knieende
  Helfer beginnt mit der 
  Herzdruckmassage. Währenddessen macht der am Kopf
  knieende Helfer die Beatmungshilfen und danach den AED
  einsatzbereit. Sofern es nicht zu einer
  Zeitverzögerung kommt, kann ein zusätzlicher Helfer der 
  Rettungsleitstelle mitteilen, dass der Patient
  reanimiert werden muss.
  \item Nach dem Betätigen des Knopfes
  \enquote{\EE{Analyse}} am AED beurteilt das Gerät, ob
  ein schockbarer Rhythmus vorliegt. Für den
  Analysevorgang muss die Herzdurckmassage unterbrochen werden. 
  Durch
  Sprachkommandos teilt der AED dem Helfer mit, ob ein Schock
  ausgelöst werden kann.
  \item Wird vom AED \EE{kein Schock empfohlen}, so ist die
  HDM unverzüglich wieder aufzunehmen. Nach zwei Minuten
  (\VonBis{5}{6} HDM-/Beatmungs-Zyklen) wird erneut eine Analyse
  durchgeführt.
  \item Wird vom AED ein \EE{Schock empfohlen}, müssen alle
  Helfer den Kontakt zum Patienten lösen und Sauerstoffquellen 
  in einem ausreichenden Sicherheitsabstand (ca. 1\Meter*) positionieren. 
  Der Helfer am
  Kopfende des Patienten bedient den AED und erinnert die
  Kollegen durch das \EE{\underline{laute}} Kommando
  \Speech{\enquote{Vorsicht! Patient nicht berühren! ---
  Schock!}} an den nötigen Sicherheitsabstand. Bevor  der
  Elektroschock durch Druck auf die Schock-Taste am AED
  endgültig ausgelöst wird, muss sich der Helfer
  vergewissen, dass niemand den Patient berührt!

  Nach der Schockabgabe wird die
  \EE{HDM sofort
  wieder aufgenommen}. Die Kombination 30 HDM und zwei
  Beatmungen wird rund zwei Minuten (\VonBis{5}{6} HDM-/Beatmungs-Zyklen)
  fortgesetzt, danach erfolgt die nächste Rhythmusanalyse.
  \item Das Vorgehen wiederholt sich jetzt bis der Patient
  Lebenszeichen zeigt oder Abbruchkriterien erfüllt sind. 
  \item Um der Erschöpfung der Helfer und damit einem
  Qualitätsverlust der HDM vorzubeugen, \EE{wird nach jedem
  HDM-/Beat\-mungs-Zyklus die Helferposition gewechselt!}
\end{itemize}
Eine grafische Übersicht über dieses Vorgehen ist in
\FullPrettyRef{tab:algorithmus-bls-aed}, dargestellt.

}{%
siehe \prettyref{tab:algorithmus-bls-aed} \ldots{}
Algorithmus BLS mit AED}{}{}{}

\subsection{Advanced Life Support (ALS)}
\TextSlide{Vorgehensweise ALS}{
Der \EE{Advanced Life Support} (ALS) beschreibt das Vorgehen
für speziell geschulte Notfallteams, welche aufgrund ihrer
Ausbildung und ihrer Ausrüstung erweiterte Maßnahmen
(venöser Zugang, Applikation von
Medikamenten, \ldots) durchführen können. Er ist daher
\E{nicht} Inhalt dieses Kapitels.
\VARIANT{VAR-STANDARD}{Dem ALS ist das Kapitel \ref{chp:ALS} gewidmet.}
}{ 
\begin{itemize}
  \item Erweiterte Maßnahmen für speziell ausgebildete
  Notfallteams
  \VARIANT{VAR-STANDARD}{\item siehe \prettyref{chp:ALS}}
\end{itemize}
}{}{}{}


\subsection{Pediatric Life Support (PLS)}
\TextSlide{Vorgehensweise PLS}{%
Die Reanimation bei Kindern wird als \I{Pediatric Life
Support} (\I{PLS}) bezeichnet. 
% Die hier vorgestellte Variante
% geht von den Basismaßnahmen (und nicht von den erweiterten
% Maßnahmen) aus und richtet sich -- so wie das BLS -- an
% Laien und Helfer aus \E{nicht} speziell ausgebildeten
% Notfallteams.
Während bei Erwachsenen der Großteil der Ursachen für einen
Herz-Kreislauf-Stillstand vom Herzen aus geht, liegt bei den
Kindern meist ein Problem des Atemwegs bzw. der Atmung (z.\,B.
Atemwegsverlegung, Ertrinken) vor. Deshalb steht -- im Gegensatz zum BLS -- beim PLS
das Freimachen der Atemwege im Vordergrund. Das
Vorgehen gemäß ERC wird wie folgt beschrieben:
\begin{itemize}
  \item \EE{Szeneüberblick mit (Selbst-)Schutz}: wie bei BLS
  \item \EE{Eindruck/Hauptbeschwerde}: wie BLS, zusätzlich
  Abschätzen des Alters
  \item \EE{Ansprechbarkeit prüfen, Bewusstsein beurteilen}:
  wie BLS, \E{kindgerecht} ansprechen
  \item Ist der \EE{Patient ansprechbar} folgt das Vorgehen
  den allgemeinen Regeln des Patientenmanagements.
  \item Ist der \EE{Patient nicht ansprechbar} wird ein
  Notarzt angefordert. Idealerweise deligiert man diese
  Aufgabe an einen zusätzlichen Helfer.
  \item \EE{Atemweg kontrollieren}: Es wird die
  Mundhöhle inspiziert und etwaige Fremdkörper
  entfernt.
  \item Beim Säugling (1.
  Lebensjahr) wird der Kopf in Neutralposition\footnote{Beim
  auf dem Rücken liegenden Säugling ist der Kopf i.\,d.\,R.
  nach vorne gebeugt, sodass eine \E{leichte} Streckung
  erforderlich sein kann um den Kopf in Neutralposition zu
  bringen \cite{Erc2010Sek06De}.} gehalten, beim Kind (Ende
  des 1. Lebensjahrs bis zur Pubertät) wird der Kopf
  vorsichtig überstreckt um die Atemwege zu öffnen. 
  \item
  \EE{Atmung prüfen}: Wie beim BLS wird für die Dauer von
  10\Sec* die Atmung des Patienten durch \E{Sehen},
  \E{Höhren} und \E{Fühlen} beurteilt. 
  \item \EE{Beurteilung} ob die \EE{Atmung normal} ist: wie
  BLS
  \item Liegt \EE{normale Atmung} vor, wird der Patient in
  die stabile Seitenlage gedreht. Danach folgen alle
  weiteren {\TxMassVit}.
  \item Danach werden \E{5 Initialbeatmungen} verabreicht.
  \item Es wird erneut nach \EE{Lebenszeichen} gesucht. Wenn
  das Kind in irgend einer Form (z.\,B.
  Spontanbewegung, Husten oder normale Atmung)
  auf die Initialbeatmungen reagiert hat, wird eine
  neuerliche Atemkontrolle durchgeführt um zu prüfen ob die
  Eigenatmung suffizient ist. Das Suchen nach Lebenszeichen
  darf nicht länger als 10\Sec* dauern. Ein gut
  ausgebildeter, erfahrener Helfer kann während dieser
  10\Sec* parallel versuchen den Puls zu tasten. Beim
  Säugling muss der Puls mehr als 60 Schläge pro Minute
  betragen.
  \item Bei \EE{Lebenszeichen}, aber
  insuffizienter Atmung, muss das Kind beatmet werden. Ist
  die Atmung suffizient wird das Kind gemäß den {\TxMassVit}
  weiter behandelt.
  \item
  Sind \EE{keine Lebenszeichen} feststellbar muss mit der
  \EE{Herz-Lung\-en-Wied\-er\-be\-leb\-ung} begonnen werden.
  \begin{itemize}
    \item Es
    wechseln \E{15 HDM} mit \E{zwei Beatmungen} ab.
    \item
    Die HDM
    erfolgt je nach Körpergröße mit zwei Fingern (Säugling)
    bzw. mit einer oder beiden Händen bei
    Kindern.\footnote{\Z{Die ERC
    empfiehlt beim Säugling alternativ auch die
    thoraxumschließende Zweidaumentechnik
    \cite{Erc2010Sek06De}.}}
    \item
    Der \EE{Druckpunkt} liegt auf der unteren Hälfte des Brustbeins. 
    \begin{itemize}
      \item Bei \E{Säuglingen} und \E{Neugeborenen} kann er dadurch
      gefunden werden, indem man den Rippenbogen bis zum Schwertfortsatz
      entlang fährt, der Druckpunkt liegt dann eine Fingerbreite darüber.
    \end{itemize}
    \item
    Die
    \E{Eindrücktiefe} soll ungefähr \nicefrac{1}{3} des
    Brustkorbs betragen. Das sind beim Säugling ca. 4\CM* und
    bei Kindern ca. 5\CM* \cite{Erc2010Sek06De}.
    \item
    Die \E{Druckfrequenz}
    (\enquote{Arbeitsgeschwindigkeit}) soll
    100-120\PerMin*{} betragen. Bei der Beatmung richtet sich
    das Atemzugvolumen ebenfalls nach der Körpergröße des
    Patienten.
    \item \Z{Gemäß ERC kann eine
    Defibrillation bei Kindern bereits ab dem 1. Lebensjahr
    erfolgen. Nachdem die meisten verwendeten
    Geräte jedoch von den Herstellern erst ab 8 Jahren
    freigegeben sind, wird bei den meisten Rettungsdiensten
    der Defibrillator nur bei Kindern über 8 Jahren
    eingesetzt. Die Vorgehensweise ist in diesem Fall gleich
    wie bei Erwachsenen.}
  \end{itemize}
\end{itemize}
Eine grafische Übersicht über dieses Vorgehen ist ein
\prettyref{tab:algorithmus-pls} dargestellt. \cite{Erc2010Sec6}

}{ 
Siehe \prettyref{tab:algorithmus-pls} \ldots{}
Algorithmus PLS 
}{}{}{}

\subsection{Reanimation Neugeborener}
\TextSlide{Reanimation des Neugeborenen}
{%
Die Reanimation des Neugeborenen wird unter \FullPrettyRef{topic:reanimation-neugeborenes} ausgeführt.}
{%
\begin{itemize}
  \item \FullPrettyRef{topic:reanimation-neugeborenes}
\end{itemize}
}

\clearpage
\begin{figure}[H]
\begin{WidePar}
\begin{center}
\Captionof{table}{Basic Life Support mit AED.\label{tab:algorithmus-bls-aed}}
\index{Basic Life Support}\index{BLS|see{Basic Life Support}}

\includegraphics[scale=1]{./images/emhofer-josef-aass/BLS}
\end{center}
\end{WidePar}
\end{figure}

%  \clearpage
\begin{figure}[H]
\begin{WidePar}
% \begin{center}
\Captionof{table}{Pediatric Life Support.\label{tab:algorithmus-pls}}

\includegraphics[scale=1]{./images/emhofer-josef-aass/PLS.pdf}\index{Pediatric Life
Support}\index{BLS|see{Pediatric Life Support}}
% \end{center}
\end{WidePar}
\end{figure}

\section{\HeadingKeyword{Ende} der Reanimation und Weiterbehandlung}

\TextSlide{Ende der Reanimation und Weiterbehandlung}
{\index{Reanimation!Beendigung}
Der Helfer darf die Wiederbelebung beenden, wenn einer der
drei angeführten Umstände (Abbruchkriterien) eintritt:

\begin{itemize}
  \item Eintreffen von \EE{weiterer Hilfe}: Der Helfer
  kann bei Eintreffen weiterer Hilfe durch diese abgelöst
  werden. I.\,d.\,R. trifft jedoch ein Notarztteam ein, mit
  welchem die Reanimation zunächst gemeinsam weitergeführt
  wird. Ein etwaiger Abbruch kann nur auf Anweisung des Notarztes
  erfolgen.
  \item \EE{Wiedererlangung des Kreislaufes} 
  (\T{ROSC}\footnote{\I{ROSC}: Return Of Spontaneous Circulation}):  
  Zeigt der Patient \EE{eindeutige Lebenszeichen} (Erwachen,
  Husten, Atmen) so wird die Reanimation vorläufig beendet. 
  Es ist sofort eine \EE{Atemkontrolle}
  durchzuführen um sicherzustellen, dass die
  Lebenszeichen richtig erkannt wurden. 
  Ist eine normale Atmung vorhanden,
  wird die Reanimation endgültig beendet.\footnote{Diese Form der Beendigung der Reanimation
  ist leider äußerst selten.} Es ist der Einschätzungsblock 
  (\FullPrettyRef{topic:einschaetzungsblock}) sowie die
  {\TxMassVit} durchzuführen. Die Sauerstoffgabe hat dabei
  einen Sp\ce{O2}-Zielwert von \VonBis{94}{98}\PC*.
  \EE{Erbrechen zählt \E{nicht} als Lebenszeichen} da dies
  durch fehlerhafte Beatmung (zu hohes Atemzugsvolumen)
  ausgelöst werden könnte. In diesem Fall wird die
  Reanimation umgehend fortgesetzt.
  \item \EE{Erschöpfung}: Sollte ein Ersthelfer all seine Kräfte
  aufgebraucht haben, wird gezwungener Maßen die
  Wiederbelebung ab- bzw.
  unterbrochen. Medizinisches Fachpersonal wird nicht in
  eine derartige Situation kommen, da weitere Hilfe
  eintreffen sollte bevor der Erschöpfungszustand eintritt.
  Außerdem ist der Erschöpfung durch Abwechseln bei der
  Durchführung der HDM und der Beatmung vorzubeugen. Der
  Wechsel findet idealerweise nach jedem
  HDM-/Beatmungszyklus, d.\,h. ca.
  alle 2\Min*, statt und darf zu keiner Verzögerung der HDM
  führen.
\end{itemize}

}{
\begin{itemize}
  \item Eintreffen von weiterer Hilfe
  \item Patient zeigt Lebenszeichen
  \begin{itemize}
    \item Erbrechen ist \E{kein} Lebenszeichen
    \item {\TxMassVit}
    \item ausführliche Untersuchung und (Fremd-)anamnese
    \end{itemize}
  \item (Erschöpfung)
\end{itemize}
}



\section{\HeadingKeyword{Anmerkungen}}
\label{topic:anmerkungen-bls}

\Text{Infektionen}{% 
Die Ansteckungsgefahr mit Krankheiten im Rahmen einer
Reanimation ist sehr gering. Es wurden noch keine
Infektionen mit schwerwiegenden Erkrankungen wie HIV oder
Hepatitis C während einer Basisreanimation beschrieben.
\cite[p.1284]{Erc2010Sec2} }{}{}{}{}

\Text{Probleme bei der HDM während des Defi-Ladevorganges}{% 
In den ERC-Leitlinien von
2010 wird erstmals empfohlen, die Herzdruckmassage fortzusetzen,
\E{während der Defibrillator sich auflädt}, d.\,h. nach der
Analyse wird weiter massiert, bis der Defi bereit zur
Schockabgabe ist. Dies soll die HDM-Pausen verringern. Bei
vielen AED-Geräten kann es jedoch zu
Problemen kommen, sie interpretieren die
reanimationsbedingten Störungen der EKG-Ableitung als
Kontaktverlust bzw. Bewegungsartefakt und brechen den
Ladevorgang ab.  \hfill
\Commentary[Herzmassage während AED-Ladezyklus]{Lt.
vorläufiger Auskunft der Distributoren ist dies momentan
zumindest mit den AED-Geräten der Fa. Schiller
(Fred{\texttrademark}, Defigard{\texttrademark} 1002) nicht
möglich, da auch während der Aufladephase das EKG über die
Defi-Elektroden überwacht wird. Die elektrische Aktivität
durch Thoraxkompressionen wird dann als Bewegungsartefakt
interpretiert und es wird eine interne Sicherheitsentladung
eingeleitet. Für die Modelle der Fred-Reihe läuft momentan
eine Evaluierung bezüglich eines Softwareupdates, es wird
daher aller Voraussicht nach in Zukunft eine Lösung geben.
 
Für die Geräte der Corpuls{\texttrademark}-Reihe liegen
uns derzeit noch keine Informationen vor.  
Der LifePak{\texttrademark} 500 von Medtronic verhält sich
wie der Defigard{\texttrademark} 1002. Bezüglich der Typen
LifePak{\texttrademark} 1000, 12 und 15 liegen ebanfalls
noch keine Informationen vor. 
Es muss
 davon ausgegangen werden, dass auch andere Geräte
ein derartiges Verhalten zeigen. \E{(Angaben ohne
Gewähr. Stand: 2011-01-19)}}

\Customized{
\begin{description}
  \item[\CustAsboe] Bei AEDs darf während der Aufladephase
  keine Herzdruckmassage durchgeführt werden. 
\end{description}
}
}{}{}{}{}

\ChapterEnd